\begin{table}[ht]
\centering
\caption[Definición de las fases del ciclo de vida]{Definición de las fases del ciclo de vida en la base de datos utilizada \citep{coutodesouza2024new}.}
\label{tab:fases_def}
\begin{tabular*}{\textwidth}{@{\extracolsep{\fill}}lp{10cm}}
\hline
Fase & Descripción Dinámica \\ 
\hline
Incipiente (Ic) & Etapa inicial de organización del vórtice, previa a la caída abrupta de presión. \\ 
% \hline
Intensificación (It) & Período de rápido crecimiento de la vorticidad ciclónica (tasa de cambio negativa intensa). \\ 
% \hline
Madurez (M) & Etapa donde el sistema alcanza su intensidad máxima (pico de vorticidad) y estabilización relativa. \\ 
% \hline
Decaimiento (D) & Fase final de debilitamiento progresivo del vórtice (ciclólisis) y disipación. \\ 
\hline
\end{tabular*}
\end{table}

% \begin{table}[ht]
% \centering
% \caption{Definición operativa de las fases del ciclo de vida en la base de datos utilizada \citep{coutodesouza2024new}.}
% \label{tab:fases_def}
% \begin{tabular*}{\textwidth}{@{\extracolsep{\fill}}lp{10cm}}
% \hline
% Fase & Descripción Dinámica \\ 
% \hline
% Incipiente (Ic) & Etapa inicial de organización del vórtice, previa a la caída abrupta de presión. \\ 
% % \hline
% Intensificación (It) & Período de rápido crecimiento de la vorticidad ciclónica (tasa de cambio negativa intensa). \\ 
% % \hline
% Madurez (M) & Etapa donde el sistema alcanza su intensidad máxima (pico de vorticidad) y estabilización relativa. \\ 
% % \hline
% Decaimiento (D) & Fase final de debilitamiento progresivo del vórtice (ciclólisis) y disipación. \\ 
% \hline
% \end{tabular*}
% \end{table}

% \begin{table}[ht]
% \centering
% \caption{Definición operativa de las fases del ciclo de vida en la base de datos utilizada \citep{coutodesouza2024new}.}
% \label{tab:fases_def}
% \begin{tabular}{l|p{10cm}}
% \hline
% \textbf{Fase} & \textbf{Descripción Dinámica} \\ \hline
% \textbf{Incipiente (Ic)} & Etapa inicial de organización del vórtice, previa a la caída abrupta de presión. \\ \hline
% \textbf{Intensificación (It)} & Período de rápido crecimiento de la vorticidad ciclónica (tasa de cambio negativa intensa). \\ \hline
% \textbf{Madurez (M)} & Etapa donde el sistema alcanza su intensidad máxima (pico de vorticidad) y estabilización relativa. \\ \hline
% \textbf{Decaimiento (D)} & Fase final de debilitamiento progresivo del vórtice (ciclólisis) y disipación. \\ \hline
% \end{tabular}
% \end{table}
